\documentclass{article}
\usepackage[margin=1in]{geometry}
\usepackage{amsmath}
\usepackage{array}
\usepackage{longtable}
\usepackage{booktabs}
\usepackage{hyperref}

\begin{document}

\section*{Detailed Match / Non-Match Map for \texttt{diag1\_compact.tex}}

This note compares every displayed equation in
\texttt{amc/examples/sample_output/diag1\_compact.tex} against the
code-backed equations described in \texttt{learn/factorized\_latex.tex}.

The goal here is not only to record matches, but also to distinguish
three cases:
\begin{itemize}
\item \textbf{Exact / explicit match}: the same equation family is named explicitly in
  \texttt{factorized\_latex.tex}, usually with a code-backed ArXiv tag.
\item \textbf{Family / packed match}: the compact equation is implemented by the code,
  but in \texttt{factorized\_latex.tex} it only appears as part of a packed matrix
  product, a grouped family, or a recoupled channel.
\item \textbf{No standalone match found}: no separate displayed equation with the same
  object name appears in \texttt{factorized\_latex.tex}; only a generic transform or
  surrounding family can be identified.
\end{itemize}

All code references below are taken from the explicit source links already
embedded in \texttt{learn/factorized\_latex.tex}.

\section*{Primary Factorized Blocks and Code Ranges}

\begin{itemize}
\item One-body intermediate, diagram I:
  \texttt{FactorizedDoubleCommutator.cc L99--L249}
  \newline
  \texttt{factorized\_latex.tex} lines 263--320.
\item One-body intermediate, diagrams IIIa and IIIb:
  \texttt{FactorizedDoubleCommutator.cc L257--L371}
  \newline
  \texttt{factorized\_latex.tex} lines 321--359.
\item Two-body intermediate, diagrams IIb and IId:
  \texttt{FactorizedDoubleCommutator.cc L380--L541}
  \newline
  \texttt{factorized\_latex.tex} lines 361--408.
\item Pandya transform for diagrams IIa and IIc:
  \texttt{FactorizedDoubleCommutator.cc L546--L801}
  \newline
  \texttt{factorized\_latex.tex} lines 410--460.
\item One-body intermediates for Ia, Ib, IVa, IVb:
  \texttt{FactorizedDoubleCommutator.cc L850--L1144}
  \newline
  \texttt{factorized\_latex.tex} lines 727--800.
\item Two-body intermediates for diagrams II and III:
  \texttt{FactorizedDoubleCommutator.cc L1151--L3049}
  \newline
  \texttt{factorized\_latex.tex} lines 801--1008.
\end{itemize}

\section*{Full Inventory}

\renewcommand{\arraystretch}{1.2}
\begin{longtable}{p{0.18\textwidth} p{0.13\textwidth} p{0.25\textwidth} p{0.16\textwidth} p{0.22\textwidth}}
\toprule
Compact equation & Status & Factorized section & Code reference & Notes \\
\midrule
\endhead

$\chi^{\alpha}$ & Exact / explicit match & One-body intermediate, diagram I & L99--L249 & Explicitly identified as the code-backed version of $\chi^{\alpha}_{ij}$; code writes the packed contraction as $\chi^{(221a)}_{de}$. \\

$f^{(I)}$ & Exact / explicit match & One-body intermediate, diagram I & L99--L249 & Explicitly identified as the code-backed version of $f^{\mathrm{I}}_{ij}$; final contraction written as $Z^I_{pq}$. \\

$\chi^{\beta}$ & Exact / explicit match & One-body intermediate, diagrams IIIa and IIIb & L257--L371 & Explicitly identified as the code-backed version of $\chi^{\beta}_{ij}$; code writes the packed intermediate as $\chi^{(221b)}_{de}$. \\

$f^{(II)}$ & Exact / explicit match & One-body intermediate, diagrams IIIa and IIIb & L257--L371 & Explicitly identified as the code-backed version of $f^{\mathrm{II}}_{ij}$. This is one of the naming swaps noted by you: the subsection is labeled by diagrams IIIa/IIIb, but the resulting compact family is $f^{(II)}$. \\

$\bar\Omega$ & No standalone match found & Closest object is the Pandya transform block for diagrams IIa and IIc & L546--L801 & The compact file has a separate displayed equation for $\bar\Omega_{ajkb}$. In \texttt{factorized\_latex.tex} there is no standalone displayed $\bar\Omega$ equation; only the generic Pandya transform $\bar X^J_{ab,cd}$ is given. \\

$\bar\chi^{\gamma}$ & Exact / explicit match & Pandya transform for diagrams IIa and IIc & L546--L801 & Explicitly identified as the code-backed version of $\bar\chi^{\gamma J}_{i\bar j k\bar l}$. The code writes the packed matrix as $\bar\chi^{(222a),J}$. \\

$f^{(III_a)}$ & Exact / explicit match & Pandya transform for diagrams IIa and IIc & L546--L801 & Explicitly identified as the code-backed version of $f^{\mathrm{III}_a}_{ij}$; written as the two contractions $Z^{IIa}_{pq}$ and $Z^{IIc}_{pq}$. \\

$\chi^{\delta}$ & Exact / explicit match & Two-body intermediate, diagrams IIb and IId & L380--L541 & Explicitly identified as the code-backed version of $\chi^{\delta}_{ijkl}$; the packed code object is $\chi^{(222b),J}$. \\

$f^{(III_b)}$ & Exact / explicit match & Two-body intermediate, diagrams IIb and IId & L380--L541 & Explicitly identified as the code-backed version of $f^{\mathrm{III}_b}_{ij}$; the final contraction is $Z^{IIb+IId}_{pq}$. \\

$\chi^{\epsilon}$ & Family / packed match & One-body intermediates for Ia, Ib, IVa, IVb & L850--L1144 & \texttt{factorized\_latex.tex} says this block packages $\chi^{\epsilon}_{ij}$ and $\chi^{\zeta}_{ij}$ together. The displayed code-backed formulas use $\chi^I_{pq}$ and $\chi^{II}_{pq}$ rather than the exact compact symbols. \\

$\chi^{\zeta}$ & Family / packed match & One-body intermediates for Ia, Ib, IVa, IVb & L850--L1144 & Same situation as $\chi^{\epsilon}$: present as a named family, but only through the packed one-body intermediates $\chi^I$ and $\chi^{II}$. \\

$\bar\chi^{\eta}$ & Family / packed match & Two-body intermediates for diagrams II and III & L1151--L3049 & Present only at family level. The text says the ordinary-channel $\chi^{\eta}$ family corresponds to $\Gamma^{\mathrm{III}_aJ}$ and the cross-coupled / double-bar family corresponds to $\overline{\overline\Gamma}^{\mathrm{III}_bJ}$. No single displayed equation with the exact compact symbol $\bar\chi^{\eta}$ is given. \\

$\chi^{\theta}$ & Family / packed match & Two-body intermediates for diagrams II and III & L1151--L3049 & Present through the tagged family ``packed intermediate for $\bar\chi^{\theta J}$'' and the later $\bar\chi_\Gamma = \bar\chi^{IV}\bar\Gamma$ construction. No one-line standalone compact formula is reproduced. \\

$\bar\chi^{\iota}$ & Family / packed match & Two-body intermediates for diagrams II and III & L1151--L3049 & Present through the tagged family ``packed intermediate for $\bar\chi^{\iota J}$ and $\chi^{\lambda J}$''. This is a family-level correspondence, not a standalone one-line reproduction. \\

$\overline{\overline\chi}^{\kappa}$ & Family / packed match & Two-body intermediates for diagrams II and III & L1151--L3049 & The factorized note says the code-packed ordinary-channel realization $-\eta\chi^{VI}-\chi^{VI,II}\eta$ corresponds to the $\chi^{\kappa}$ family, i.e. to $\Gamma^{\mathrm{IV}_aJ}$. The exact double-bar compact symbol is not displayed as its own standalone equation. \\

$\chi^{\lambda}$ & Family / packed match & Two-body intermediates for diagrams II and III & L1151--L3049 & Present through the tagged family ``packed intermediate for $\bar\chi^{\iota J}$ and $\chi^{\lambda J}$''. The exact expanded compact form is not printed separately in the factorized note. \\

$\Gamma^{(I)}$ & Family / packed match & One-body intermediates for Ia, Ib, IVa, IVb & L850--L1144 & Explicit family-level match. The factorized note says this block packages $\Gamma^{\mathrm{I}J}_{ijkl}$ and $\Gamma^{\mathrm{II}J}_{ijkl}$ together in a single update formula $Z^J_{pq,rs}$. \\

$\Gamma^{(II)}$ & Family / packed match & One-body intermediates for Ia, Ib, IVa, IVb & L850--L1144 & Explicit family-level match, but not a one-equation-per-family layout. The compact equation corresponds to the part of the combined update that uses $\chi^{II}$. \\

$\Gamma^{(III_a)}$ & Family / packed match & Two-body intermediates for diagrams II and III & L1151--L3049 & Explicitly mapped as the ordinary-channel $\chi^{\eta}$ family. The code-backed form is the packed matrix update $\chi^{III}\Gamma + h_Z\Gamma(\chi^{III})^T$, not a standalone expanded compact equation. \\

$\overline{\overline\Gamma}^{(III_b)}$ & Family / packed match & Two-body intermediates for diagrams II and III & L1151--L3049 & Explicitly mapped by the recoupled $\chi^{\eta}$ family and the inverse Pandya return to standard coupling. Again, this is a packed / recoupled realization, not a literal compact-form rewrite. \\

$\bar\Gamma^{(III_c)}$ & Family / packed match & Two-body intermediates for diagrams II and III & L1151--L3049 & Explicitly mapped through the barred $\chi^{\theta}$ family and then through the packed product $\bar\chi_\Gamma = \bar\chi^{IV}\bar\Gamma$. \\

$\Gamma^{(IV_a)}$ & Family / packed match & Two-body intermediates for diagrams II and III & L1151--L3049 & Explicitly mapped as the $\chi^{\kappa}$ family through the packed ordinary-channel update $-\eta\chi^{VI}-\chi^{VI,II}\eta$. \\

$\overline{\overline\Gamma}^{(IV_b)}$ & Family / packed match & Two-body intermediates for diagrams II and III & L1151--L3049 & Explicitly mapped as part of the barred / double-barred $\chi^{\iota}$ and $\chi^{\lambda}$ families with inverse Pandya insertion. \\

$\bar\Gamma^{(IV_c)}$ & Family / packed match & Two-body intermediates for diagrams II and III & L1151--L3049 & Explicitly mapped together with the $\chi^{\lambda}$ family; the packed inverse-Pandya forms for $\bar\Gamma^{\mathrm{III}_cJ}$ and $\bar\Gamma^{\mathrm{IV}_cJ}$ are written near the end of the section. \\

\bottomrule
\end{longtable}

\section*{What Matches Exactly}

The strongest one-to-one matches are:
\begin{itemize}
\item $\chi^{\alpha}$ and $f^{(I)}$.
\item $\chi^{\beta}$ and $f^{(II)}$.
\item $\bar\chi^{\gamma}$ and $f^{(III_a)}$.
\item $\chi^{\delta}$ and $f^{(III_b)}$.
\end{itemize}

These are the cases where \texttt{factorized\_latex.tex} explicitly says
``this is the code-backed version of'' the same named compact family and gives
an associated \texttt{FactorizedDoubleCommutator.cc} block.

\section*{What Only Matches at Family Level}

The remaining compact equations are implemented in the factorized code, but the
factorized note groups them into packed matrix products, recoupled channels, or
shared update formulas:
\begin{itemize}
\item $\chi^{\epsilon}$, $\chi^{\zeta}$.
\item $\bar\chi^{\eta}$, $\chi^{\theta}$, $\bar\chi^{\iota}$,
  $\overline{\overline\chi}^{\kappa}$, $\chi^{\lambda}$.
\item $\Gamma^{(I)}$, $\Gamma^{(II)}$, $\Gamma^{(III_a)}$,
  $\overline{\overline\Gamma}^{(III_b)}$, $\bar\Gamma^{(III_c)}$,
  $\Gamma^{(IV_a)}$, $\overline{\overline\Gamma}^{(IV_b)}$,
  $\bar\Gamma^{(IV_c)}$.
\end{itemize}

So these are \emph{matches}, but not line-by-line compact reproductions.

\section*{What Does Not Have a Standalone Match}

Only one displayed equation in \texttt{diag1\_compact.tex} clearly lacks a
standalone counterpart in \texttt{factorized\_latex.tex}:
\begin{itemize}
\item $\bar\Omega_{ajkb}$.
\end{itemize}

The factorized note clearly uses the same object implicitly through the generic
Pandya transform, but it does not print a separate displayed equation for
$\bar\Omega$ itself.

\section*{Naming Mismatches to Keep in Mind}

\begin{itemize}
\item In the one-body factorized section, the compact family $f^{(II)}$ is tied
  to the subsection called ``diagrams IIIa and IIIb''. This is a real naming /
  regrouping mismatch, not a mistake in your reading.
\item In the two-body factorized section, $\Gamma^{(I)}$ and $\Gamma^{(II)}$ are
  presented together in one block, and the same happens for several later
  $\Gamma$ families.
\item The factorized note often documents the packed code object rather than the
  final expanded compact expression. For that reason, several compact equations
  have a valid code match even though no visually identical formula is printed.
\end{itemize}

\section*{Bottom Line}

Comparing \texttt{diag1\_compact.tex} against \texttt{factorized\_latex.tex} and
using only the code references embedded there, I find:
\begin{itemize}
\item 8 explicit / near-one-to-one matches.
\item 15 family-level or packed matches.
\item 1 equation with no standalone displayed counterpart: $\bar\Omega$.
\end{itemize}

So almost everything in \texttt{diag1\_compact.tex} is represented in the
factorized implementation note, but much of the two-body material is only
visible there in grouped packed form rather than as one compact equation per
family.

\end{document}
